\section{Empirical Insights}

We summarise five findings from the \SAFK{} experiments most
relevant to practitioners deploying AI agents in security
environments. Each is paired with the operational implication it
carries for SOC teams.

\subsection{Knowledge Mismatch Drives Hallucination}

\textbf{Finding.} Across 750 LLM-based agent executions,
hallucination rate increases monotonically with knowledge mismatch
($\delta_K$): from 1.4\% at $\delta_K = 0$ to 39.1\% at $\delta_K = 1.0$
(Pearson $\rho = 0.92$, $p < 0.001$). The pattern holds across all three
LLM architectures (LangGraph, AutoGen, CrewAI). The rule-based baseline
(AUT-4) hallucinates zero times regardless of $\delta_K$, confirming
hallucination as an LLM-specific failure mode.

\textbf{Implication.} Verify threat-knowledge coverage before
autonomous execution. Knowledge gaps are the single strongest
predictor of dangerous output fabrication, which makes
retrieval-augmented generation (RAG) pipeline quality a first-order
reliability concern: the completeness of the retrieved context
directly determines hallucination risk.

\subsection{Different Mismatches Produce Distinct Failure Signatures}

\textbf{Finding.} Each mismatch dimension produces a diagnostically
separable failure pattern (chi-square test of independence,
$p < 0.001$). Conditioning on the LLM-agent records whose
\textit{primary} mismatch dimension is at level $\geq 0.5$:
capability deprivation ($\delta_C$) predominantly produces
\textit{Paralysis} (78\%); skill deficit ($\delta_S$) produces silent
\textit{Omission} (72\%); knowledge gaps ($\delta_K$) produce
\textit{Hallucination} (32\%) and \textit{Misclassification} (21\%);
reasoning perturbation ($\delta_R$) produces \textit{Confabulation}
(38\%) and \textit{Overclaiming} (17\%).

\textbf{Implication.} The manifestation pattern is a diagnostic
fingerprint. A SOC observing high paralysis should investigate tool
access and API permissions, not LLM quality. A high confabulation rate
suggests reasoning corruption or adversarial input manipulation rather
than knowledge gaps. This enables targeted remediation rather than
wholesale agent replacement.

\subsection{ARS Generalises as a Reliability Gate}

\textbf{Finding.} The learned ARS threshold ($\theta^* = 0.58$)
achieves test F1 = 0.96 across all five threat classes
(Table~\ref{tab:e3_mag}), with perfect precision (1.00) and recall
ranging from 0.84 to 1.00. The minimal train-test gap (0.95 vs.\ 0.96)
indicates that ARS captures a generalisable reliability signal rather
than overfitting to specific threat classes or agent architectures.
The reported F1 uses the post-execution ARS that includes $\delta_R$;
in deployment, the pre-execution gate uses $\delta_S$, $\delta_K$,
$\delta_C$, and $\delta_T$, with $\delta_R$ added after execution
for diagnosis and continuous improvement.

\begin{table}[t]
\caption{ARS Anomaly Detection Performance (Test Set)}
\label{tab:e3_mag}
\centering
\footnotesize
\begin{tabular}{lccc}
\toprule
\textbf{Threat Class} & \textbf{Precision} & \textbf{Recall} & \textbf{F1} \\
\midrule
Ransomware     & 1.00 & 0.94 & 0.97 \\
Phishing       & 1.00 & 0.94 & 0.97 \\
Insider Threat & 1.00 & 0.87 & 0.93 \\
DDoS           & 1.00 & 0.84 & 0.91 \\
APT Intrusion  & 1.00 & 1.00 & 1.00 \\
\midrule
\textbf{Overall} & \textbf{1.00} & \textbf{0.92} & \textbf{0.96} \\
\bottomrule
\end{tabular}
\end{table}

\textbf{Implication.} ARS works operationally as a pre-execution
deployment gate. Perfect precision means no false alarms: every
flagged execution truly exhibited anomalous behaviour, removing the
common adoption blocker where reliability tooling generates so many
false alerts that SOC teams disable it.

\subsection{Multi-Agent Pipelines Amplify Failures}

\textbf{Finding.} Comparing AUT-1 (single-agent LangGraph) with
AUT-2 (three-stage AutoGen pipeline) shows a 5.7\% headline ARS gap
that conflates framework identity with pipeline depth. A controlled
ablation on 27 strictly matched mismatch buckets isolates the two
contributions: framework identity contributes \textbf{0.13\%}
(essentially zero) and pure pipeline depth contributes
\textbf{17.25\%} relative ARS reduction, consistent across all five
threat classes.

\textbf{Implication.} Decomposing a security task into a multi-agent
pipeline does not improve reliability; it reduces it whenever
component agents carry mismatch vulnerabilities. Prefer well-validated
single-agent designs for critical tasks. If multi-agent is required,
add inter-stage verification checkpoints and confidence-gated
propagation so that low-confidence intermediate outputs do not silently
seed downstream stages.

\subsection{Epistemic Gap Separates Knowledge Risk from Confidence Risk}

\textbf{Finding.} The Epistemic Gap ($\mathcal{E}$) and raw knowledge
mismatch ($\delta_K$) achieve statistically equivalent correlation
with false-negative severity ($\rho_\mathcal{E} = \rho_{\delta_K} =
0.80$, Fisher $z$-test $p = 0.78$). On 30 live CISA KEV incidents
beyond every model's training cutoff, five frontier LLMs scored
identically on surface accuracy (vendor 67\%, product 63\%, CVE
$\leq$3\%), yet stated confidence diverged sharply: GPT-5.4-mini
reported mean confidence 0.78 on 29/30 records while Claude Opus 4.7
abstained on 27/30.\footnote{Model identifiers reflect the public API
versions current at evaluation time (April 2026); \texttt{gpt-5.4-mini}
and \texttt{claude-opus-4-7} are accessed through the OpenAI and
Anthropic APIs respectively.} Same accuracy, different operational
risk.

\textbf{Implication.} The Epistemic Gap's value is not raw predictive
lift but \textit{interpretability}: it factorises risk into a
``what the agent does not know'' term ($\delta_K$) and a
``how confidently it acts under that ignorance'' term ($1 - U$),
exposing two independently controllable mitigations. Calibration is
therefore as actionable a lever as knowledge coverage. The
hallucination-versus-$\delta_K$ relationship replicates directionally
across four LLM families from 7B to frontier scale (bucket-aggregated
$\rho \in [0.78, 0.95]$), so the reliability signal is structural
rather than backbone-specific.
